% !TEX root = ../main.tex

\chapter{Design}
\label{ch:design}

\section{Top-level overview}
\label{sec:top-level}

\begin{figure}[H]
  \centering
  \resizebox{\textwidth}{!}{%
    \begin{tikzpicture}[
        font=\small,
        blk/.style={draw, rounded corners=2pt, minimum height=1.5cm,
        text width=2.4cm, align=center, fill=gray!8},
        core/.style={blk, text width=2.6cm, fill=blue!10},
        small/.style={draw, rounded corners=2pt, minimum height=0.8cm,
          text width=2.3cm, align=center, fill=gray!8,
        font=\footnotesize},
        lbl/.style={font=\footnotesize},
        arr/.style={-{Latex[length=2mm]}, thick},
        side/.style={-{Latex[length=2mm]}, densely dashed},
        ctrl/.style={-{Latex[length=1.6mm]}, densely dotted, red!70!black},
        node distance=0.8cm
      ]
      % ----- main datapath -----
      \node[blk]                 (cls) {Input\\classification};
      \node[blk,  right=of cls]  (rr)  {Range\\reduction};
      \node[core, right=of rr]   (cor) {CORDIC core\\(40 stages)};
      \node[blk,  right=of cor]  (rec) {Quadrant\\reconstruction};
      \node[blk,  right=of rec]  (nrm) {Normalization\\and rounding};

      % ----- output selector (trapezoid) -----
      \coordinate (mx) at ($(nrm.east)+(0.9,0)$);
      \draw[fill=gray!8] ($(mx)+(0,0.75)$) -- ($(mx)+(0.6,0.4)$)
      -- ($(mx)+(0.6,-1.4)$) -- ($(mx)+(0,-1.75)$) -- cycle;
      \coordinate (mxout) at ($(mx)+(0.6,-0.5)$);
      \coordinate (mxtop) at ($(mx)+(0.3,0.58)$);

      % ----- data arrows -----
      \draw[arr] ($(cls.west)+(-1.0,0)$) -- node[above, lbl]{$x$} (cls.west);
      \draw[arr] (cls) -- (rr);
      \draw[arr] (rr)  -- node[above, lbl]{$r$} (cor);
      \draw[arr] (cor) -- node[above, lbl]{$x_n, y_n$} (rec);
      \draw[arr] (rec) -- (nrm);
      \draw[arr] (nrm) -- (mx);
      \draw[arr] (mxout) -- node[above, lbl]{$\sin x$} ++(1.1,0);

      % ----- small-argument path -----
      \node[small, above=0.55cm of cor] (sma) {Small argument\\$r \cdot \pi/2$};
      \draw[arr] (rr.north) |- (sma.west);
      \draw[arr] (sma.east) -| (rec.north);

      % ----- side information -----
      \coordinate (sideY) at ($(sma.north)+(0,0.45)$);
      \draw[side] (cls.north) -- (cls.north |- sideY) -- (mxtop |- sideY)
      node[pos=0.5, above, lbl]
      {side information: sign, $k \bmod 4$, input class
      (delayed with the data)}
      -- (mxtop);

      % ----- bypass path -----
      \coordinate (byY)   at ($(cls.south)+(0,-0.45)$);
      \coordinate (byin)  at (mx |- byY);
      \coordinate (byout) at ($(cls.south)+(0.5,0)$);
      \draw[arr] (byout) -- (byout |- byY) -- (byin)
      node[pos=0.45, below, lbl]
      {bypass path (delayed): $\pm 0$, $x$ itself, NaN};

      % ----- stall control -----
      \coordinate (stY) at ($(byY)+(0,-0.95)$);
      \coordinate (rdy) at ($(mxout)+(1.1,-0.45)$);
      \node[right, lbl, text=red!70!black] at (rdy) {ready};
      \draw[densely dotted, red!70!black] (rdy) |- (cls.south |- stY)
      node[pos=0.75, below, lbl, text=red!70!black]
      {stall: all stages hold their values while the receiver is not ready};
      \foreach \b in {cls, rr, cor, rec, nrm}
      \draw[ctrl] (\b.south |- stY) -- (\b.south);
    \end{tikzpicture}%
  }
  \caption{Top-level architecture of the CORDIC-based sine unit}
  \label{fig:top-level}
\end{figure}